\documentclass[11pt]{article}

\usepackage[margin=1in]{geometry}
\usepackage[T1]{fontenc}
\usepackage[utf8]{inputenc}
\usepackage{amsmath,amssymb}
\usepackage{graphicx}
\usepackage{booktabs}
\usepackage[round]{natbib}
\usepackage{hyperref}
\hypersetup{colorlinks=true, linkcolor=blue!50!black, citecolor=blue!50!black, urlcolor=blue!50!black}
\usepackage{enumitem}
\usepackage{xcolor}
\usepackage{caption}

\title{\LARGE Progress Report\\[4pt]
\large Prompt-Driven Interactive Segmentation as a Dataset-Annotation Tool}
\author{
  \begin{tabular}{c}
    Student A: 2301498\\
    Abdelhafeez Aref Abdelhafeez Almaz
  \end{tabular}
  \and
  \begin{tabular}{c}
    Student B: 2501925\\
    Abdelrahman Eid Ashour Mohamed
  \end{tabular}
}
\date{Milestone 2 --- \today\\CSE 681 --- Digital Image Processing}

\begin{document}
\maketitle

\begin{abstract}
This project investigates whether point- and box-prompted SAM 2.1 can reduce the
effort required to produce instance masks while retaining agreement with COCO
ground truth. The local dataset workflow now contains 300 deterministically
selected COCO 2017 validation images, 979 non-crowd instances across 80
categories, and exported instance masks. The React/TypeScript interface and
FastAPI backend now support browsing those images, selecting a target instance,
and saving a prediction with its annotation time and computed IoU/Dice. The
software checks pass, but SAM 2 inference has not yet been run end-to-end and
neither method has a measured annotation-quality or timing result; no
performance conclusion is claimed in this report.
\end{abstract}

\section{Recap}
\label{sec:recap}
Manual instance-mask annotation is labor-intensive. We compare a manual polygon
baseline with a human-in-the-loop workflow in which an annotator prompts a
frozen SAM 2.1 Hiera-Tiny model and accepts or refines its mask. Both methods
are intended to annotate the same COCO object instances. The data source is
COCO 2017 val2017 \citep{lin2014coco}; the reproducible subset contains 300
images selected with seed 42 from images with 1--8 non-crowd instances. The
selected annotations yield 979 non-crowd instance masks in 80 categories. The
reference mask is intended for scoring only, not for choosing prompts.

\section{Implementation Status}
\label{sec:status}
The current application uses a React, TypeScript, and Vite frontend with a
Python FastAPI backend. The backend exposes local COCO image listing and image
serving, accepts prompted image-segmentation requests, and provides an
annotation-saving endpoint that checks image/annotation pairing, computes
binary-mask IoU and Dice, and records the mask, prompts, and elapsed time in a
local evaluation manifest. The frontend includes COCO image and target-instance
selectors, positive/negative point and box prompts, mask review, and acceptance
with metric display. Downloaded images, masks, checkpoint, and evaluation
outputs are excluded from Git; reproduction instructions are tracked.
The images retain their individual terms of use \citep{cocoTerms}.

The COCO data/API tests pass (16 backend tests in total), and the frontend
TypeScript build and lint checks pass. These are software-validation results,
not segmentation experiments. The optional SAM 2 Python dependency has not
been installed in the current environment, so the downloaded checkpoint and
inference route have not yet been validated together. The original plan's
manual polygon baseline also remains to be run. No model training is planned:
the contribution is an interactive annotation and evaluation workflow around a
frozen SAM 2 model \citep{ravi2024sam2}.

\subsection{Baseline Reproduction}
The manual polygon baseline has not yet been reproduced or measured. No
annotation times, IoU, or Dice values are available for it. It remains a
required M3 experiment: annotate the same target objects with a standard
polygon tool, using the same timer start/stop definition, then compare masks
with the corresponding COCO references.

\subsection{Your Method}
The web workflow is implemented through the software path: browse a local COCO
sample, choose one of its non-crowd objects, add points and/or a box, request a
mask, and accept it. Acceptance saves the prediction and elapsed time and
computes metrics against the matching exported COCO mask. Inference latency is
reported separately from annotation time. The remaining end-to-end gap is
execution with the actual SAM 2.1 checkpoint and collection of evaluated
annotations; until then the user interface and API are not evidence of model
quality or time savings.

\section{Preliminary Results}
\label{sec:results}
The current quantitative evidence describes dataset preparation and software
validation only. No human-created baseline masks or SAM 2 predictions have
been evaluated, so method-level IoU, Dice, and annotation-time comparisons are
not yet available.

\begin{table}[ht]
\centering
\caption{Evidence available at the time of this progress report. ``Not
measured'' denotes missing experimental data, not a zero score.}
\label{tab:prelim}
\begin{tabular}{@{}lccc@{}}
\toprule
\textbf{Workflow} & \textbf{Instances} & \textbf{Mean IoU / Dice} & \textbf{Annotation time} \\
\midrule
Manual polygon baseline & 0 evaluated & Not measured & Not measured \\
SAM 2.1 prompted tool & 0 evaluated & Not measured & Not measured \\
\midrule
Prepared COCO subset & 979 masks & -- & -- \\
Automated backend tests & 16 passed & -- & -- \\
\bottomrule
\end{tabular}
\end{table}

No qualitative SAM 2 examples are included because a real inference run has
not yet been completed. The next report should include representative
successful cases and failure cases captured from actual annotations, without
exposing the reference mask while prompts are selected.

\section{Challenges \& Deviations from the M1 Plan}
\label{sec:challenges}
The local subset and mask export are ready, but the optional model package has
not been installed and real inference has not been exercised. Consequently,
the planned manual-versus-prompted comparison is behind the software
implementation. The M2 work prioritized a reproducible data/API path and
persisted metrics rather than claiming results from synthetic test fixtures.
COCO source images and derived masks remain local and ignored by Git; the
tracked data guide describes how to reproduce them. The work therefore
deviates from a results-oriented milestone report: the implementation path is
in place, while experimental validation remains incomplete.

\section{Updated Work Breakdown \& MVP Status}
\label{sec:workbreakdown}
The software portion of the MVP is implemented, but the M1 minimum viable
experimental deliverable has not been reached: the fixed-sample manual and
prompted masks, paired timing data, and comparative metrics are still missing.
The team must complete these measurements before drawing conclusions in M3.
Because the team has two members, Student B combines the planned model/tool
and evaluation/integration responsibilities.

\begin{table}[ht]
\centering
\caption{Updated two-member milestone allocation and status.}
\label{tab:tasks}
\begin{tabular}{@{}p{0.20\linewidth}p{0.35\linewidth}p{0.35\linewidth}@{}}
\toprule
\textbf{Milestone} & \textbf{Student A} & \textbf{Student B} \\
\midrule
M2 (current) &
COCO subset and manual-baseline protocol prepared; polygon measurements remain pending. &
React/FastAPI annotation and evaluation workflow implemented; real SAM 2 inference and evaluated samples remain pending. \\
M3 (planned) &
Complete the manual polygon baseline on the fixed sample and record time consistently. &
Install and validate SAM 2.1, complete prompted annotations, audit metrics, and produce the comparison and qualitative examples. \\
\bottomrule
\end{tabular}
\end{table}

\section{Next Steps}
\label{sec:next}
\begin{enumerate}[leftmargin=*]
  \item Install the optional SAM 2 dependency, configure the downloaded
  Hiera-Tiny checkpoint, and verify an end-to-end point- and box-prompt run on
  local COCO images.
  \item Select and record the fixed evaluation sample and target annotation IDs
  before either method is run; do not reveal ground truth during prompt choice.
  \item Collect paired manual polygon and prompted masks and times using the
  same timer definition. Record device, model/checkpoint identity, prompts,
  failures, and any excluded cases.
  \item Run the evaluator, verify per-instance mask dimensions and IDs, and
  report mean/median annotation time, mean IoU and Dice, and representative
  successes and failures. Avoid time-savings claims until these measurements
  are complete.
\end{enumerate}

\bibliographystyle{unsrtnat}
\bibliography{references}

\end{document}
